\documentclass[a4paper,9pt]{extarticle}

%=========================================================
% PACKAGES
%=========================================================

\usepackage[
a4paper,
top=1.8cm,
bottom=1.5cm,
footskip=7mm,
left=1cm,
right=1cm,
headheight=14pt,
headsep=5mm
]{geometry}
\usepackage[T1]{fontenc}
\usepackage[utf8]{inputenc}
\usepackage[french]{babel}

\usepackage[sfdefault]{FiraSans}
\renewcommand{\familydefault}{\sfdefault}

\usepackage{multicol}
\setlength{\columnsep}{0.7cm}

\usepackage{amsmath,amssymb}

\usepackage{titlesec}
\usepackage{enumitem}

\usepackage{tcolorbox}
\tcbuselibrary{skins,breakable}

\usepackage{xcolor}

\usepackage{graphicx}
\usepackage{float}

\usepackage{tikz}

\usepackage{fancyhdr}
\usepackage[hidelinks]{hyperref}


%=========================================================
% PAGE SETTINGS
%=========================================================

\setlength{\parindent}{0pt}
\setlength{\parskip}{1.5pt}

\setlength{\abovedisplayskip}{4pt}
\setlength{\belowdisplayskip}{4pt}

\setlist[itemize]{
leftmargin=*,
itemsep=1pt,
topsep=1pt
}


%=========================================================
% COLORS
%=========================================================

\definecolor{Blue}{HTML}{E3F0FF}
\definecolor{Green}{HTML}{E3F7EA}
\definecolor{Pink}{HTML}{FFE5ED}
\definecolor{Orange}{HTML}{FFF0DC}
\definecolor{Purple}{HTML}{EEE7FF}
\definecolor{Yellow}{HTML}{FFFADB}


\definecolor{BlueTitle}{HTML}{285B8F}
\definecolor{GreenTitle}{HTML}{287052}
\definecolor{PinkTitle}{HTML}{A84063}
\definecolor{OrangeTitle}{HTML}{A96500}
\definecolor{PurpleTitle}{HTML}{6245A5}


%=========================================================
% HEADER
%=========================================================

\pagestyle{fancy}

\fancyhf{}

\fancyhead[L]{\small\textbf{Computer Security}}
\fancyhead[C]{\small Mandatory Access Control}
\fancyhead[R]{\small Week 04 --- 28/09-04/10/2026}

\fancyfoot[C]{\small\href{https://github.com/AlphaXT38G/notes_ba5}{github.com/AlphaXT38G/notes\_ba5}}
\fancyfoot[R]{\small \thepage}

\renewcommand{\headrulewidth}{0.4pt}



%=========================================================
% TITLES
%=========================================================

\titleformat{\section}
{\large\bfseries\color{BlueTitle}}
{}{0pt}{}

\titlespacing{\section}
{0pt}{6pt}{3pt}


% separator

\newcommand{\separator}{
\vspace{2pt}
\hrule
\vspace{3pt}
}


%=========================================================
% BOXES
%=========================================================

% DEFINITION - blue
\newtcolorbox{definition}{
enhanced,
breakable,
colback=Blue,
colframe=BlueTitle,
boxrule=0.7pt,
arc=0mm,
title=\textbf{Definition},
fonttitle=\bfseries
}

% IMPORTANT - pink
\newtcolorbox{important}{
enhanced,
breakable,
colback=Pink,
colframe=PinkTitle,
boxrule=0.7pt,
arc=0mm,
title=\textbf{Key idea},
fonttitle=\bfseries
}

% FORMULAS - green
\newtcolorbox{formula}{
enhanced,
breakable,
colback=Green,
colframe=GreenTitle,
boxrule=0.7pt,
arc=0mm,
title=\textbf{Key formulas},
fonttitle=\bfseries
}

% WARNING - orange
\newtcolorbox{warning}{
enhanced,
breakable,
colback=Orange,
colframe=OrangeTitle,
boxrule=0.7pt,
arc=0mm,
title=\textbf{Warning},
fonttitle=\bfseries
}

% EXAMPLE - purple
\newtcolorbox{example}{
enhanced,
breakable,
colback=Purple,
colframe=PurpleTitle,
boxrule=0.7pt,
arc=0mm,
title=\textbf{Example},
fonttitle=\bfseries
}

% SUMMARY - orange
\newtcolorbox{summary}{
enhanced,
breakable,
colback=Yellow,
colframe=OrangeTitle,
boxrule=0.7pt,
arc=0mm,
title=\textbf{Summary},
fonttitle=\bfseries
}



%=========================================================
% FIGURE BOX
%=========================================================

\newtcolorbox{figurebox}{
enhanced,
colback=white,
colframe=gray,
boxrule=0.5pt,
arc=0mm,
center
}



%=========================================================
% DOCUMENT
%=========================================================

\begin{document}

%\footnotesize   --> To reduce the size of the text

\thispagestyle{fancy}



\begin{multicols}{2}


%=========================================================
\section{Mandatory Access Control}


\begin{definition}

\textbf{Mandatory Access Control}

Access to and operations on resources are determined by the security policy

\begin{itemize}

\item ``owner'' may not exist or not have power to set permissions against policy

\item the security policy must be enforced despite subjects trying to subvert it

\end{itemize}

\end{definition}


\begin{definition}

\textbf{SECURITY MODEL:} a design pattern for a specific security
property or set of properties

\end{definition}



%=========================================================
\section{Bell-La Padula (BLP) model: Protecting confidentiality}


Subjects S and objects O associated to a level of confidentiality

\textbf{Subjects access rights are defined by four attributes:}

\begin{itemize}

\item \textbf{Execute:} the subject cannot see or modify the object, but can run it

\item \textbf{Read:} the subject can only see the object but cannot modify it

\item \textbf{Append:} the subject cannot read the object, but can add attach new content

\item \textbf{Write:} the subject can see the object and add content or modify existing content

\end{itemize}

These access rights are defined in an access control matrix


\separator


\textbf{Level function for objects: Classification}

Objects are associated to a Security Level

(they have a label, and belong to one or more categories)


\begin{formula}

Security Level = (Classification, \{set of categories\})

\begin{itemize}

\item \textbf{Classification} -- total order of labels (e.g., Unclassified, Confidential, Secret, Top Secret)

\item \textbf{Categories} -- compartments of objects with a common topic (e.g., Nuclear, NATO, Crypto)

\end{itemize}

\end{formula}


\begin{definition}

\textbf{DOMINANCE RELATIONSHIP}

A security level $(l_1,c_1)$ ``dominates'' $(l_2, c_2)$ if and only if $l_1 \geq l_2$ and $c_2$ is a subset of $c_1$

\end{definition}


\begin{important}

\textbf{Three key properties}

(a) Dominates is transitive.

(b) Top and bottom elements.

(c) Only partial order

\end{important}


\separator


\textbf{Level function for subjects: Clearance level}

\begin{itemize}

\item \textbf{Clearance} -- maximum security level a subject has been assigned: clearance level($S_i$)

\item \textbf{Current security level} -- subjects can operate at lower security levels: current-level($S_i$)

\end{itemize}


\begin{definition}

\textbf{SIMPLE SECURITY PROPERTY (SS-PROPERTY)}

If (subject, object, r) is a current access, then level(subject) dominates level(object)

\end{definition}


\begin{definition}

\textbf{STAR PROPERTY (*-PROPERTY)}

if a subject has simultaneous ``observe'' (r,w) access to $O_1$
and ``alter'' (a,w) access to $O_2$
then level($O_2$) dominates level($O_1$)

\end{definition}


\begin{definition}

\textbf{DISCRETIONARY PROPERTY (DS-PROPERTY)}

if an access (subject, object, action) takes place it must be in the access
control matrix

\end{definition}



%=========================================================
\section{BLP: Basic Security Theorem}


\begin{important}

\textbf{BASIC SECURITY THEOREM}

if all state transitions are security preserving, and the initial state is
secure, then every subsequent state is secure regardless of the inputs

\end{important}


If for any individual access:

(1) the ss-property is preserved

(2) the *-property is preserved

(3) the ds-property is preserved

\dots{} then for any sequential composition security holds!


\separator


\begin{definition}

\textbf{COVERT CHANNEL}

any channel that allows information flows contrary to the security policy

\end{definition}


\begin{definition}

\textbf{DECLASSIFICATION}

remove classification label

\end{definition}



%=========================================================
\section{The BIBA model for integrity}


\begin{important}

\textbf{Two key rules (strict)}

\begin{itemize}

\item \textbf{simple integrity (no read-down):} protects higher integrity principals from being corrupted
by lower integrity level data

\item \textbf{*-integrity (no write-up):} prevents lower integrity principals from corrupting high integrity
data

\end{itemize}

\end{important}


\textbf{Low-water-mark policy for subjects}

\begin{itemize}

\item Subjects start processes at their highest integrity level.

\item When accessing an object, its current level is lowered to
the lowest of the two: current-level(s) and level(o)

\end{itemize}

Temporary downgrade for the session

\begin{itemize}

\item Example: mitigate impact of a network Trojan

\item Hard to avoid label creep

\end{itemize}


\begin{definition}

\textbf{SANITIZATION}

Process of taking objects with ``low'' integrity and ``lifting them'' to ``high integrity''

\end{definition}


\begin{important}

\textbf{PRINCIPLE 2: FAIL-SAFE DEFAULT}

``Base access decisions on permission rather than exclusion'' [SS75]

\end{important}



%=========================================================
\section{Chinese Wall model: Access rules}


\begin{important}

A subject can read an object (for either read or write) if the access does not
allow an information flow between items with labels in the same conflict set

\end{important}



\end{multicols}


\end{document}
